\documentclass[../../../include/open-logic-section]{subfiles}

\begin{document}

\olfileid{sth}{ordinals}{basic}
\olsection{క్రమసంఖ్యల ప్రాథమిక ధర్మాలు}

తొలి కొన్ని క్రమసంఖ్యలే
సహజ సంఖ్యలని చూశాం.
క్రమసంఖ్యల సిద్ధాంతాన్ని
అభివృద్ధి చేయడానికి
ముఖ్య కారణం సహజ
సంఖ్యలపై వర్తించే
ఆగమన సూత్రాన్ని
విస్తరించడం. ప్రాథమిక
ఫలితాల క్రమం ద్వారా
దానికి చేరుకుంటాం.

\begin{lem}\ollabel{ordmemberord}
క్రమసంఖ్యలోని ప్రతి
\tetoken{మూలకం}{!!{element}}
క్రమసంఖ్యే.
\end{lem}

\begin{proof}
$\alpha$ క్రమసంఖ్య,
$b \in \alpha$ అనుకుందాం.
$\alpha$ సంక్రామకం
కాబట్టి $b \subseteq \alpha$.
అందువల్ల $\in$,
$\alpha$ను సుక్రమితం
చేసినట్టే $\in$, $b$నూ
సుక్రమితం చేస్తుంది.

$b$ సంక్రామకమని
చూడడానికి $x \in c \in b$
అనుకుందాం. $b
\subseteq \alpha$ కాబట్టి
$c \in \alpha$.
మళ్లీ $\alpha$ సంక్రామకం
కాబట్టి $c \subseteq
\alpha$; అందువల్ల
$x \in \alpha$.
కాబట్టి $x, c, b \in \alpha$.
అయితే $\in$, $\alpha$ను
సుక్రమితం చేస్తుంది;
కనుక \olref[sth][ordinals][wo]{wo:strictorder}
ప్రకారం $\alpha$పై
$\in$ సంక్రామక సంబంధం.
అందువల్ల $x
\in c \in b$ నుంచి
$x \in b$ వస్తుంది.
సాధారణీకరిస్తే,
$c \subseteq b$.
\end{proof}

\begin{cor}\ollabel{ordissetofsmallerord}
ఏ క్రమసంఖ్య~$\alpha$కైనా
$\alpha = \Setabs{\beta \in \alpha}{\beta \text{ ఒక క్రమసంఖ్య}}$.
\end{cor}

\begin{proof}
\olref{ordmemberord}
నుంచి వెంటనే వస్తుంది.
\end{proof}

తదుపరి రెండు ప్రధాన
ఫలితాలు,
\olref{ordinductionschema},
\olref{ordtrichotomy},
స్థూలంగా చెప్పేది:
క్రమసంఖ్యలన్నీ సభ్యత్వం
ద్వారా సుక్రమితమైనవి.

\begin{thm}[అతిపరిమిత ఆగమనం]\ollabel{ordinductionschema}
ఏ $\phi(x)$ సూత్రానికైనా:
\[
	\text{ఒకవేళ }\exists \alpha \phi(\alpha)\text{ అయితే }\exists \alpha(\phi(\alpha)
	\land  (\forall \beta \in \alpha) \lnot \phi(\beta))
\]
ఇక్కడ చూపిన పరిమాణకారకాలు
క్రమసంఖ్యలకు మాత్రమే
పరిమితమని అంతర్లీనంగా
అనుకుంటాం.
\end{thm}

\begin{proof}
%\olref{ordinductionschema1}.
ఏదో క్రమసంఖ్య
$\alpha$కు $\phi(\alpha)$
అనుకుందాం. $ (\forall \beta
\in \alpha) \lnot \phi(\beta)$
అయితే నిరూపణ పూర్తయింది.
లేకపోతే $\alpha$
క్రమసంఖ్య కాబట్టి,
దానిలో సూత్రం $\phi$ను
సంతృప్తిపరిచే మూలకాలలో
$\in$-అత్యల్ప
\tetoken{మూలకం}{!!{element}}
ఉంటుంది; అది
\olref{ordmemberord}
ప్రకారం క్రమసంఖ్యే.
ఆ మూలకానికి ముందున్న
ఏ క్రమసంఖ్యా ఆ
సూత్రాన్ని సంతృప్తిపరచదు.
\sourcecorrection{OLTESTORDBASIC-001}{మూలం అత్యల్ప మూలకం ``\(\phi\)'' అని మాత్రమే చెబుతుంది; ఇక్కడ \(\phi\)ను సంతృప్తిపరిచే మూలకాలలో అత్యల్పమని స్పష్టం చేశాం.}
%	\olref{ordinductionschema2}. ఏదో క్రమసంఖ్య
%	$\gamma$ ఉంది, దానికి $\lnot\phi(\gamma)$. అప్పుడు
%	\olref{ordinductionschema1} ప్రకారం $\in$-కనిష్ఠ క్రమసంఖ్య
%	$\alpha$కు $\lnot\phi(\alpha)$. కాబట్టి $(\forall \beta <
%	\alpha) \phi(\beta)$; దీని వల్ల సోపాధికపు పూర్వభాగం
%	అసత్యం అవుతుంది.
\end{proof}
\noindent
\olref{ordinductionschema}ను
ఈ కింది పథకంగా కూడా
చెప్పవచ్చని గమనించండి:
\[
\text{ఒకవేళ }\forall \alpha((\forall \beta \in \alpha)\phi(\beta) \lif 
\phi(\alpha))\text{ అయితే }\forall \alpha\phi(\alpha)
\]
దీనికి \olref{ordinductionschema}లో
$\lnot\phi(\alpha)$ను
తీసుకొని, ప్రాథమిక
తార్కిక మార్పులు చేస్తే
చాలు.

\begin{thm}[త్రివిధత]\ollabel{ordtrichotomy}
ఏ క్రమసంఖ్యలు
$\alpha$, $\beta$కైనా
$\alpha \in \beta \lor \alpha = \beta \lor \beta \in \alpha$.
\end{thm}

\begin{proof}
నిరూపణ ద్వంద్వ ఆగమనం
ద్వారా; అంటే
\olref{ordinductionschema}ను
రెండుసార్లు వాడతాం.
$x \in y \lor x = y \lor y \in x$
అయితే $x$, $y$తో
\emph{పోల్చదగినది}
అంటాం.

ఆగమనానికి, $\alpha$లోని
ప్రతి క్రమసంఖ్యా
\emph{ప్రతి} క్రమసంఖ్యతో
పోల్చదగినదని అనుకుందాం.
మరొక ఆగమనానికి,
$\alpha$, $\beta$లోని
ప్రతి క్రమసంఖ్యతో
పోల్చదగినదని అనుకుందాం.
$\alpha$, $\beta$తో
పోల్చదగినదని చూపుతాం.
అప్పుడు $\beta$పై
ఆగమనం ద్వారా $\alpha$
ప్రతి క్రమసంఖ్యతో
పోల్చదగినదని వస్తుంది;
తరువాత $\alpha$పై
ఆగమనం ద్వారా
\emph{ప్రతి} క్రమసంఖ్యా
\emph{ప్రతి} క్రమసంఖ్యతో
పోల్చదగినదని వస్తుంది.
కావలసింది ఇదే.
$\alpha \notin \beta$,
$\beta \notin
\alpha$ అనుకొని
$\alpha = \beta$
అని చూపితే సరిపోతుంది.

$\alpha \subseteq \beta$
అని చూపడానికి,
$\gamma \in \alpha$ను
స్థిరపరచండి; అది
\olref{ordmemberord}
ప్రకారం క్రమసంఖ్య.
కాబట్టి మొదటి ఆగమన
పరికల్పన ప్రకారం,
$\gamma$, $\beta$తో
పోల్చదగినది. కానీ
$\gamma
= \beta$ అయినా,
$\beta \in \gamma$
అయినా $\beta \in \alpha$
వస్తుంది ($\alpha$
సంక్రామకమనే విషయాన్ని
అవసరమైతే వాడుతూ);
అది మన ఊహకు
విరుద్ధం. కాబట్టి
$\gamma \in \beta$.
సాధారణీకరిస్తే,
$\alpha \subseteq
\beta$.

రెండవ ఆగమన పరికల్పనతో
అచ్చం ఇలాంటి వాదన
$\beta \subseteq \alpha$
అని చూపుతుంది.
కాబట్టి $\alpha = \beta$.
\end{proof}\noindent
అందువల్ల కొన్నిసార్లు
$\alpha \in \beta$ బదులు
$\alpha <\beta$ అని
రాస్తాం; ఎందుకంటే
$\in$ క్రమ సంబంధంలా
ప్రవర్తిస్తుంది. పరిచయం,
అలాగే $\alpha \in
\beta \lor \alpha = \beta$
కంటే $\alpha \leq \beta$
రాయడం సులభం కావడం
తప్ప దీనికి లోతైన
కారణం లేదు.\footnote{$\alpha
\mathrel{\underline{\in}} \beta$
అని రాయవచ్చు; కానీ
అది పూర్తిగా సంప్రదాయేతరం.}

చివరి ఫలితాలకు రెండు
సత్వర పర్యవసానాలు
ఇవి; మొదటిది మన
కొత్త సంకేతలిపిని
ఉపయోగిస్తుంది:

\begin{cor}\ollabel{ordordered}
$\exists \alpha\phi(\alpha)$
అయితే
$\exists \alpha(\phi(\alpha)
\land \forall \beta(\phi(\beta) \lif \alpha \leq \beta))$.
ఇంకా ఏ క్రమసంఖ్యలు
$\alpha, \beta, \gamma$కైనా,
$\alpha
\notin \alpha$ మరియు
$\alpha \in \beta \in \gamma \lif \alpha \in
\gamma$ రెండూ ఉంటాయి.
\end{cor}

\begin{proof}
\olref[wo]{wo:strictorder}
మాదిరిగానే.
\end{proof}

\begin{prob}
\olref[sth][ordinals][basic]{ordtrichotomy}
నిరూపణలోని ``అచ్చం
ఇలాంటి వాదన''ను
పూర్తి చేయండి.
\end{prob}

\begin{cor}\ollabel{corordtransitiveord}
$A$ క్రమసంఖ్య అప్పుడే
మరియు అప్పుడే $A$
క్రమసంఖ్యల సంక్రామక
సమితి.
\end{cor}

\begin{proof}
\emph{ఎడమ నుంచి కుడికి.}
\olref{ordmemberord}
ప్రకారం. \emph{కుడి నుంచి
ఎడమకి.} $A$ క్రమసంఖ్యల
సంక్రామక సమితి అయితే,
\olref{ordinductionschema},
\olref{ordtrichotomy}
ప్రకారం $\in$, $A$ను
సుక్రమితం చేస్తుంది.
\end{proof}

\olref{ordinductionschema},
\olref{ordtrichotomy} ప్రకారం
$\in$ క్రమసంఖ్యలను
సుక్రమితం చేస్తుందని
స్థూలంగా చెప్పాం.
అయితే ఈ కింది ఫలితం
వల్ల అలాంటి వాదనలతో
\emph{చాలా జాగ్రత్తగా}
ఉండాలి:

\begin{thm}[బురాలి-ఫోర్టీ వైరుధ్యం]\ollabel{buraliforti}
అన్ని క్రమసంఖ్యల
సమితి లేదు.
\end{thm}

\begin{proof}
వైరుధ్యాన్ని చూపడానికి,
$O$ అన్ని క్రమసంఖ్యల
సమితి అనుకుందాం.
$\alpha \in
\beta \in O$ అయితే,
\olref{ordmemberord}
ప్రకారం $\alpha$
క్రమసంఖ్య; కాబట్టి
$\alpha \in O$.
అందువల్ల $O$
సంక్రామకం; అందుచేత
\olref{corordtransitiveord}
ప్రకారం $O$
క్రమసంఖ్య. అప్పుడు
$O \in O$; ఇది
\olref{ordordered}కు
విరుద్ధం.
\end{proof}
\noindent
ఈ ఫలితం
\citeauthor{Burali-Forti1897}
పేరుతో పిలవబడుతుంది.
అయితే అన్ని క్రమసంఖ్యల
సమితి ఉందనే ఊహలోని
\emph{వైరుధ్యాన్ని}
స్పష్టంగా మొదట గుర్తించింది
కాంటర్; 1899లో
డెడెకిండ్‌కు రాసిన
ఉత్తరంలో. వాన్
హేయెనూర్ట్ వివరించినట్టు:
\begin{quote}
  బురాలి-ఫోర్టీ స్వయంగా
  వైరుధ్య నిరూపణ ద్వారా,
  క్రమసంఖ్యల సహజ క్రమం
  కేవలం పాక్షిక క్రమమే
  అని స్థాపించబడిందని
  భావించాడు.
  \citep[p.~105]{Heijenoort1967}
\end{quote}
బురాలి-ఫోర్టీ తప్పును
పక్కనపెట్టి, పై విషయాన్ని
ఇలా సంక్షిప్తంగా చెప్పవచ్చు:
క్రమసంఖ్యలు సభ్యత్వం
ద్వారా ఒక్కొక్కటిగా
సుక్రమితమైన సమితులు;
అవి మొత్తంగా ఒక సమితిని
ఏర్పరచకపోయినా,
సమష్టిగా సభ్యత్వం
ద్వారా సుక్రమితమైనవి.

ఇక్కడ ముగింపుగా,
\olref{ordinductionschema},
\olref{ordtrichotomy}
నుంచి వచ్చే క్రమసంఖ్యల
మరిన్ని ప్రాథమిక
ధర్మాలు ఇవి.

\begin{prop}
క్రమసంఖ్యల ఏ కఠినంగా
దిగే అనుక్రమమైనా
పరిమితమైనదే.
\end{prop}

\begin{proof}
అనంతమైన కఠినంగా దిగే
క్రమసంఖ్యల అనుక్రమం
$\alpha_0 > \alpha_1 > \alpha_2 > \ldots$
లో $<$-కనిష్ఠ
మూలకం ఉండదు;
ఇది \olref{ordinductionschema}కు
విరుద్ధం.
\end{proof}

\begin{prop}\ollabel{ordinalsaresubsets}
ఏ క్రమసంఖ్యలు
$\alpha, \beta$కైనా
$\alpha \subseteq \beta \lor \beta \subseteq \alpha$.
\end{prop}

\begin{proof}
$\alpha \in \beta$ అయితే,
$\beta$ సంక్రామకం
కాబట్టి $\alpha \subseteq \beta$.
అలాగే $\beta \in \alpha$
అయితే $\beta \subseteq
\alpha$. $\alpha = \beta$
అయితే $\alpha \subseteq \beta$
మరియు $\beta \subseteq \alpha$.
కాబట్టి \olref{ordtrichotomy}
ప్రకారం ఫలితం వస్తుంది.
\end{proof}

\begin{prop}\ollabel{ordisoidentity}
ఏ క్రమసంఖ్యలు
$\alpha, \beta$కైనా
$\alpha = \beta$ అప్పుడే
మరియు అప్పుడే
$\ordeq{\alpha}{\beta}$.
\end{prop}

\begin{proof}
క్రమసంఖ్యలు సుక్రమాలు;
కాబట్టి త్రివిధత
(\olref[basic]{ordtrichotomy}),
\olref[iso]{wellordnotinitial}
నుంచి వెంటనే వస్తుంది.
\end{proof}

\begin{prob}\ollabel{probunionordinalsordinal}
$X$లోని ప్రతి మూలకమూ
క్రమసంఖ్య అయితే,
$\bigcup X$ కూడా
క్రమసంఖ్య అని నిరూపించండి.
\end{prob}

\end{document}
